\chapter{Proporcionalidad y datos}\label{ch:g6:propdata}

Si tres cuadernos cuestan $6$ euros, cuestan seis cuadernos $12$: duplicar el
Cuadernos, el doble de precio. Las cantidades que se comportan así son
\emph{\omterm{def:g6:propdata:def}{proporcional}} --- una de las ideas más útiles de todos
matemáticas, desarrolladas aún más en \cref{ch:g7:prop}. El capítulo termina
con lectura y dibujo de tablas de datos simples.

\section{Cantidades proporcionales}

\begin{definition}[Proporcionalidad]\label{def:g6:propdata:def}
Dos cantidades son \emph{proporcionales}\index{proporcionales} cuando el
Los valores de uno se obtienen a partir de los valores del otro multiplicando.
siempre por el \emph{mismo} número, llamado el \emph{proporcionalidad
coeficiente}.
\end{definition}

\begin{example}\label{ex:g6:propdata:table}
Cuadernos en $2$ euros cada uno:
\begin{center}
\renewcommand{\arraystretch}{1.3}
\begin{tabular}{l|cccc}
notebooks & $3$ & $5$ & $8$ & $12$ \\
\hline
price (euros) & $6$ & $10$ & $16$ & $24$ \\
\end{tabular}
\end{center}
Cada precio es el número de cuadernos. $\times 2$: el coeficiente es
$2$ (el precio unitario). Una comprobación rápida de que una mesa está \omterm{def:g6:propdata:def}{proporcional}:
toda la ``columna \omterm{def:g3:division:remainder}{cocientes}'' $\frac{6}{3}, \frac{10}{5}, \frac{16}{8},
\frac{24}{12}$ deben ser iguales --- aquí todos son iguales $2$.
\end{example}

\begin{example}[Un no ejemplo]\label{ex:g6:propdata:nonexample}
La edad y la altura no son \omterm{def:g6:propdata:def}{proporcional}: a $12$-años no es el doble
alto como un $6$-años de edad. Consulta los números: $1.50$ m en $12$ años y
$1.15$ m en $6$ años dan \omterm{def:g3:division:remainder}{cocientes} $\frac{1.50}{12} = 0.125$ y
$\frac{1.15}{6} \approx 0.19$ --- no es igual.
\end{example}

\begin{method}[Completar una tabla de proporcionalidad]\label{met:g6:propdata:complete}
Tres formas, a elegir libremente:
\begin{enumerate}
  \item \emph{coeficiente}: encuentra el multiplicador de uno completo
        columna, aplicarla a las demás;
  \item \emph{columnas}: una columna se puede multiplicar por un número o dos
        se pueden agregar columnas para producir una nueva columna;
  \item \emph{volver a la unidad}: encuentra el valor de \emph{uno} artículo
        Primero, luego multiplica.
\end{enumerate}
\end{method}

\begin{example}\label{ex:g6:propdata:methods}
Cinco tazas idénticas cuestan $15$ euros; ¿Cuánto cuestan ocho tazas?

\emph{Volver a la unidad}: una taza cuesta $15 \div 5 = 3$ euros, entonces ocho
costo de las tazas $8 \times 3 = 24$ euros.

\emph{columnas}: $8 = 5 + 3$; tres tazas cuestan $\frac{3}{5}$ de $15$,
es decir.\ $9$; entonces ocho tazas cuestan $15 + 9 = 24$ euros. La misma respuesta, ya que
debe ser.
\end{example}

\section{Porcentajes}

\begin{definition}[Porcentaje]\label{def:g6:propdata:percent}
``$ t\,\%$ de una cantidad'' significa la \omterm{def:g6:fractions:def}{fracción} $\frac{t}{100}$ de ello:
$25\,\%$ de $60$ es $\frac{25}{100} \times 60 = 15$. Un porcentaje es un
\omterm{def:g6:propdata:def}{proporcionalidad} con coeficiente $\frac{t}{100}$.
\end{definition}

\begin{example}\label{ex:g6:propdata:percent}
una clase de $30$ los estudiantes contienen $40\,\%$ chicas. Número de niñas, paso
por paso: $40\,\%$ medio $\frac{40}{100} = 0.4$, y
$0.4 \times 30 = 12$ chicas. Puntos de referencia útiles:
$50\,\%$ es uno \omterm{def:g3:division:half}{medio}, $25\,\%$ uno \omterm{def:g3:division:half}{cuarto}, $10\,\%$ una décima parte ---
entonces $10\,\%$ de $30$ es $3$, y $40\,\% = 4 \times 10\,\%$ es
$4 \times 3 = 12$: misma respuesta, calculada mentalmente.
\end{example}

\section{Lectura y dibujo de gráficos.}

\begin{method}[Leer una tabla o un gráfico]\label{met:g6:propdata:read}
Cualquiera que sea la imagen: tabla, gráfico de barras, gráfico de líneas:
\begin{enumerate}
  \item lea primero el título y las unidades en ambos ejes;
  \item para responder una pregunta, ubique la barra/columna derecha y luego lea
        el valor cuidadosamente contra la graduación;
  \item para comparar, compare barras por altura --- pero verifique el eje
        comienza en $0$, o la imagen puede inducir a error.
\end{enumerate}
\end{method}

\begin{example}\label{ex:g6:propdata:chart}
El siguiente gráfico de barras muestra la cantidad de libros prestados en la escuela.
biblioteca durante una semana.

\begin{omfigure}
\begin{tikzpicture}
\begin{axis}[
  width=8.6cm, height=5.2cm,
  ybar, bar width=16pt,
  symbolic x coords={Mon,Tue,Wed,Thu,Fri},
  xtick=data,
  ymin=0, ymax=22,
  ytick={5,10,15,20},
  ylabel={books borrowed},
  tick label style={font=\small},
  label style={font=\small}]
  \addplot[fill=omDef!30, draw=omDef] coordinates
    {(Mon,12) (Tue,8) (Wed,17) (Thu,6) (Fri,15)};
\end{axis}
\end{tikzpicture}

{\small Un gráfico de barras: una barra por día, altura igual al recuento.}
\end{omfigure}

Lecturas: el día más ocupado es el miércoles ($17$ libros). martes y
Jueves juntos representan $8 + 6 = 14$ libros --- menos que el viernes
solo ($15$). Total de la semana:
$12 + 8 + 17 + 6 + 15 = 58$ libros.
\end{example}

\begin{omfigure}
\begin{tikzpicture}
\begin{axis}[omaxis,
  width=8.2cm, height=5.4cm,
  xmin=0, xmax=9.4, ymin=0, ymax=28,
  xlabel={notebooks}, ylabel={price (euros)},
  xtick={1,2,3,4,5,6,7,8,9}, ytick={4,8,12,16,20,24}]
  \addplot[omDef, thick, domain=0:9] {3*x};
  \addplot[omDef, only marks, mark=*, mark size=1.5pt] coordinates
    {(1,3) (2,6) (3,9) (4,12) (5,15) (6,18) (7,21) (8,24)};
  \draw[dashed, omThm] (axis cs:6,0) -- (axis cs:6,18)
    -- (axis cs:0,18);
\end{axis}
\end{tikzpicture}

{\small \omterm{def:g6:propdata:def}{Proporcional} cantidades tienen una gráfica muy reconocible: la
Los puntos se alinean en una línea recta que pasa por el origen (aquí cuadernos en
$3$ euros cada uno; los guiones leen el precio de $6$ cuadernos: $18$
euros).}
\end{omfigure}

\section{Ejercicios}

\begin{exercise}[$\star $]\label{exo:g6:propdata:1}
es la mesa \omterm{def:g6:propdata:def}{proporcional}? Justificar computacionalmente \omterm{def:g3:division:remainder}{cocientes}.
\begin{center}
\renewcommand{\arraystretch}{1.2}
\begin{tabular}{l|ccc}
kg of apples & $2$ & $3$ & $5$ \\
\hline
price (euros) & $5$ & $7.5$ & $12.5$ \\
\end{tabular}
\qquad
\begin{tabular}{l|ccc}
age (years) & $2$ & $4$ & $8$ \\
\hline
height (cm) & $85$ & $103$ & $130$ \\
\end{tabular}
\end{center}
\end{exercise}

\begin{exercise}[$\star $]\label{exo:g6:propdata:2}
Cuatro croissants cuestan $4.80$ euros. Encuentra el precio de un croissant,
luego de siete croissants.
\end{exercise}

\begin{exercise}[$\star $]\label{exo:g6:propdata:3}
Completa el \omterm{def:g6:propdata:def}{proporcionalidad} tabla (¡coeficiente primero!):
\begin{center}
\renewcommand{\arraystretch}{1.2}
\begin{tabular}{l|cccc}
liters of fuel & $10$ & $25$ & $?$ & $60$ \\
\hline
price (euros) & $18$ & $?$ & $72$ & $?$ \\
\end{tabular}
\end{center}
\end{exercise}

\begin{exercise}[$\star $]\label{exo:g6:propdata:4}
un auto usa $6$ L de combustible por $100$ km. ¿Cuánto combustible para $250$ kilómetros? Para
$350$ kilómetros? ¿Hasta dónde puede llegar con $27$ ¿L?
\end{exercise}

\begin{exercise}[$\star $]\label{exo:g6:propdata:5}
Calcula mentalmente, utilizando los puntos de referencia de \cref{ex:g6:propdata:percent}:
$50\,\%$ de $84$; \quad $25\,\%$ de $200$; \quad $10\,\%$ de $63$;
\quad $30\,\%$ de $70$.
\end{exercise}

\begin{exercise}[$\star $]\label{exo:g6:propdata:6}
en una escuela de $450$ estudiantes, $60\,\%$ comer en la cafetería. cuantos
estudiantes es eso? ¿Cuantos no lo hacen?
\end{exercise}

\begin{exercise}[$\star $]\label{exo:g6:propdata:7}
Usando el gráfico de barras de \cref{ex:g6:propdata:chart}: en qué días fueron
menos que $10$ libros prestados? ¿Cuántos libros más se tomaron prestados?
¿Miércoles que jueves?
\end{exercise}

\begin{exercise}[$\star $]\label{exo:g6:propdata:8}
Las temperaturas al mediodía de lunes a viernes fueron $14$, $16$, $13$,
$17$, $20$ grados. Dibuje un gráfico de barras de estos datos (elija una opción sensata).
graduación) y lea el día más cálido.
\end{exercise}

\begin{exercise}[$\star\star $]\label{exo:g6:propdata:9}
una receta para $6$ la gente necesita $450$ g de harina y $3$ huevos. Adáptalo
para $10$ gente. (Para los huevos, piensa antes de escribir un \omterm{def:g6:decimals:places}{número decimal}
de huevos!)
\end{exercise}

\begin{exercise}[$\star\star $]\label{exo:g6:propdata:10}
A velocidad constante, un ciclista viaja $24$ km en $1$ hora.
\begin{enumerate}
  \item ¿Hasta dónde viaja? $2$ h? En \omterm{def:g3:division:half}{medio} una hora? En
        $1$ h $30$ min?
  \item ¿Cuánto tiempo necesita para $60$ kilómetros?
\end{enumerate}
\end{exercise}

\begin{exercise}[$\star\star $]\label{exo:g6:propdata:11}
Una tienda ofrece ``$20\,\%$ fuera de todo ''. Calcule el descuento y
el nuevo precio de: una pelota en $15$ euros; un escándalo en $40$ euros. es el
nuevo precio \omterm{def:g6:propdata:def}{proporcional} al precio anterior? ¿Cuál es el coeficiente?
\end{exercise}

\begin{exercise}[$\star\star\star $]\label{exo:g6:propdata:12}
Se encienden dos velas de la misma altura al mismo tiempo. el grueso
se quema completamente en $6$ horas, el delgado en $4$ horas, cada una a las
su propia tasa constante. ¿Después de cuánto tiempo exactamente la vela delgada?
\omterm{def:g3:division:half}{medio} ¿Tan alto como el grueso? (Exprese las alturas restantes después
$ t $ horas como \omterm{def:g6:fractions:def}{fracciones} de la altura inicial.)
\end{exercise}

\section{Problema: mapas, planos y cómo mentir con una tabla}

\begin{problem}[{Problema de fin de semana --- La escala de un mapa es una proporcionalidad: las longitudes siguen el coeficiente, las áreas siguen su cuadrado y los gráficos mal dibujados engañan a la vista.}]
\label{pb:g6:propdata:1}
Cada mapa, plano y modelo obedece a una regla: longitudes reales y
las longitudes dibujadas son \emph{\omterm{def:g6:propdata:def}{proporcional}}
(\cref{def:g6:propdata:def}). El coeficiente tiene un nombre famoso.
--- el \emph{escalar}\index{escalar} --- y esconde dos trampas que
este problema surge deliberadamente: \omterm{def:g6:measure:area}{áreas} hacer \emph{no} sigue el
escala, y los porcentajes, como los gráficos, dependen enteramente de lo que
se miden contra.

\medskip
\noindent\textbf{Parte I --- Reading a map.}
Un mapa de senderismo a escala $1:100\,000$: cada centímetro del mapa
representa $100\,000$ centímetros sobre el suelo.
\begin{enumerate}
  \item Convertir $100\,000$ cm en kilómetros. ¿Qué hace? $1$ centímetros
        de este mapa representan, en km?
  \item Dos pueblos mienten $7.5$ cm de distancia en el mapa; la orilla de
        un lago es un $12$ cm de curva. Da las distancias reales.
  \item Una calzada romana totalmente recta discurre $20$ km. ¿Cuánto dura?
        en el mapa?
  \item Un segundo mapa anuncia ``$1$ centímetros para $5$ kilómetros''. Escribe su
        escala en la forma $1: \,?$. ¿Cuál de los dos mapas muestra?
        ¿Más detalles para la misma región?
  \item Haz una mesa pequeña (distancia en el mapa). $1$, $2$, $7.5$ centímetros
        contra la distancia real) para el mapa de caminatas, y explique
        por qué una escala es exactamente una \omterm{def:g6:propdata:def}{proporcionalidad} en el sentido de
        \cref{def:g6:propdata:def}. ¿Cuál es el coeficiente que
        ¿Convierte centímetros en el mapa en kilómetros?
\end{enumerate}

\medskip
\noindent\textbf{Parte II --- Zoe's bedroom plan.}
Zoe dibuja un plano de su dormitorio: un rectángulo de $4$ m por
$3$ m --- a escala $1:50$.
\begin{enumerate}[resume]
  \item ¿Cuáles son las dimensiones de la habitación en el plano?
  \item Su cama mide $2$ m por $0.9$ m, y la puerta es
        $80$ centímetros de ancho. Da las tres medidas en el plano.
  \item Ahora la trampa. calcular lo real \omterm{def:g6:measure:area}{área} de la habitación en
        m $^2$, entonces el \omterm{def:g6:measure:area}{área} de su plano en cm $^2$. Convertir el
        real \omterm{def:g6:measure:area}{área} en centímetros $^2$
        (\cref{def:g6:measure:area}) y dividir por el plan
        \omterm{def:g6:measure:area}{área}: es el \omterm{def:g3:division:remainder}{cociente} $50$?
  \item Explica el número que encontraste: en un $1:50$ plan, cada
        Un centímetro cuadrado de papel representa un cuadrado real de
        $50$ centímetros por $50$ centímetro. cuantos cm reales $^2$ es eso? Estado
        la regla: cuando las longitudes se dividen por $50$, \omterm{def:g6:measure:area}{áreas} son
        dividido por\,\dots
  \item A \omterm{def:g4:numbers:round}{redondo} cubiertas de alfombras $6$ centímetro $^2$ en el plano. que real
        \omterm{def:g6:measure:area}{área} cubre, en cm $^2$ y luego en m $^2$?
\end{enumerate}

\medskip
\noindent\textbf{Parte III --- How to lie with a chart (and with
a percentage).}
\begin{enumerate}[resume]
  \item una tienda vendida $105$ el valor de euros el sábado y $110$ en
        Domingo. El gerente dibuja un gráfico de barras cuya vertical
        El eje comienza en $100$: el listón del sábado sube $5$ pequeño
        unidades, el bar dominical $10$. ¿Qué significa la imagen?
        sugieren sobre el domingo, y ¿cuál es la verdad? (Calcular
        El aumento del domingo como porcentaje del sábado, al
        porcentaje más cercano.) ¿Qué advertencia de
        \cref{met:g6:propdata:read} ¿El gerente lo ignoró?
  \item Vuelva a dibujar (o describa) el gráfico honesto, con el eje
        comenzando en $0$: ¿Cómo se comparan las dos barras ahora?
  \item Una colección crece a partir de $40$ a $60$ sellos: calcular el
        aumentar como porcentaje del valor inicial. entonces
        retrocede ante $60$ a $40$: calcula la disminución como
        un porcentaje de \emph{es} valor inicial. ¿Por qué son los
        dos respuestas diferentes, para lo mismo $20$ ¿sellos?
  \item Temporada de rebajas: un abrigo en $100$ euros obtiene ``$20\,\%$
        apagado'', y en la caja un ``extra $10\,\%$ fuera del
        precio reducido''. Calcule el precio final paso a paso.
        es el descuento total $30\,\%$? Explíquele al comprador
        lo que realmente es.
  \item Finale, de vuelta en el mapa de la pregunta 4 ($1$ centímetros para
        $5$ km): un bosque ocupa $3$ centímetro $^2$ de ese mapa. que
        es real \omterm{def:g6:measure:area}{área}, en kilómetros $^2$? Concluimos con la regla de
        todo este problema: las longitudes siguen la escala, \omterm{def:g6:measure:area}{áreas}
        sigue su\,\dots
\end{enumerate}
\end{problem}
